\documentclass[11pt,a4paper]{article}

\usepackage[utf8]{inputenc}
\usepackage[indonesian]{babel}
\usepackage{amsmath, amssymb, amsfonts}
\usepackage{geometry}
\geometry{margin=1in, headheight=14pt}
\usepackage{graphicx}
\usepackage{booktabs}
\usepackage{hyperref}
\usepackage{tcolorbox}
\tcbuselibrary{skins,breakable}
\usepackage{listings}
\usepackage{xcolor}
\usepackage{fancyhdr}

\definecolor{unibblue}{RGB}{0, 51, 102}
\definecolor{unibgold}{RGB}{212, 175, 55}
\definecolor{darkgreen}{RGB}{34, 139, 34}

\hypersetup{
    colorlinks=true,
    linkcolor=unibblue,
    citecolor=darkgreen,
    urlcolor=blue
}

\pagestyle{fancy}
\fancyhf{}
\rhead{\textbf{TKE-41XX: Optimisasi untuk Teknik Elektro}}
\lhead{Catatan Kuliah Minggu ke-2}
\rfoot{Halaman \thepage}

\lstset{
  language=Python,
  basicstyle=\ttfamily\footnotesize,
  keywordstyle=\color{unibblue}\bfseries,
  commentstyle=\color{darkgreen}\itshape,
  stringstyle=\color{red!70!black},
  numbers=left,
  numberstyle=\tiny\color{gray},
  numbersep=5pt,
  breaklines=true,
  frame=single,
  rulecolor=\color{unibblue!30},
  tabsize=4
}

\newcommand{\R}{\mathbb{R}}

\title{\textbf{Catatan Kuliah Minggu ke-2:}\\Optimisasi Linier (LP) dan Metode Grafis 2D}
\author{\textbf{Ir. Novalio Daratha S.T., M.Sc., Ph.D.}\\Program Studi S1 Teknik Elektro --- Universitas Bengkulu}
\date{Semester Ganjil 2026}

\begin{document}

\maketitle

\begin{tcolorbox}[colback=unibblue!5,colframe=unibblue,title=\textbf{Ringkasan Eksekutif Pertemuan ke-2}]
Dokumen ini membahas dasar-dasar matematis dan geometris dari *Linear Programming* (LP). Mahasiswa mempelajari cara pembentukan ruang solusi poligon konveks feasibel, pencarian titik pojok (*corner point feasible solutions*), serta solusi komputasional menggunakan `scipy.optimize.linprog`.
\end{tcolorbox}

\tableofcontents
\newpage

\section{Asumsi dan Karakteristik Utama Linear Programming}
Masalah *Linear Programming* (LP) memiliki empat asumsi utama:
\begin{enumerate}
    \item **Proposionalitas**: Kontribusi setiap variabel keputusan terhadap fungsi tujuan dan kendala bersifat sebanding secara linier.
    \item **Aditivitas**: Total kontribusi fungsi tujuan dan kendala merupakan penjumlahan dari kontribusi masing-masing variabel secara terpisah.
    \item **Divisibilitas**: Variabel keputusan dapat bernilai pecahan/kontinu ($x_j \in \R$).
    \item **Kepastian (Deterministik)**: Seluruh koefisien $c_j, a_{ij}, b_i$ diketahui dengan pasti.
\end{enumerate}

\section{Geometri Ruang Solusi 2D dan Metode Grafis}
Untuk masalah LP dua variabel keputusan ($x_1, x_2$), ruang feasibel dapat digambarkan secara 2D:
\begin{enumerate}
    \item Gambarkan garis batas untuk setiap kendala pertidaksamaan $a_{i1} x_1 + a_{i2} x_2 = b_i$.
    \item Tentukan daerah feasibel yang memenuhi seluruh pertidaksamaan sekaligus ($A \mathbf{x} \le \mathbf{b}, \mathbf{x} \ge 0$).
    \item Plot garis tingkat fungsi tujuan $Z = c_1 x_1 + c_2 x_2$ (garis isoprofit/isocost).
    \item Geser garis $Z$ searah pertumbuhan nilai tujuan hingga menyentuh titik pojok (*corner point*) terakhir dari poligon feasibel.
\end{enumerate}

\section{Implementasi Program Python}

\begin{lstlisting}[caption={Skrip Python Solusi LP scipy.optimize.linprog}]
import numpy as np
from scipy.optimize import linprog

# Maksimalkan Z = 500 x1 + 800 x2 => minimalkan -500 x1 - 800 x2
c = [-500, -800]
A = [[2, 1], [1, 3]]
b = [40, 90]

res = linprog(c, A_ub=A, b_ub=b, bounds=(0, None), method='highs')

print(f"x1 Optimal = {res.x[0]:.2f}")
print(f"x2 Optimal = {res.x[1]:.2f}")
print(f"Keuntungan Maksimum = Rp {-res.fun:.2f}")
\end{lstlisting}

\section{Kesimpulan}
1. Metode grafis memberikan pemahaman intuisi geometris yang sangat kuat mengenai sifat titik pojok ekstrim (*corner point*) pada masalah LP.
2. Pemrograman LP dengan SciPy HighS solver memberikan kecepatan komputasi yang sangat tinggi dan presisi presisi.

\end{document}
